%\documentclass{beamer}
%\usepacakge[UTF8]{ctex}
%----------------------------------------------------------------------------------------
%    PACKAGES AND THEMES
%----------------------------------------------------------------------------------------

\documentclass[aspectratio=169,xcolor=dvipsnames]{beamer}
\usepackage[UTF8]{ctex}
\usetheme{SimplePlus}
\setbeamertemplate{footline}{}

\usepackage{hyperref}
\usepackage{graphicx} % Allows including images
\usepackage{booktabs} % Allows the use of \toprule, \midrule and \bottomrule in tables

\usepackage{mathpartir}

\usepackage{tikz-cd}

\usepackage{comment}

\usepackage{fontspec}
\setmonofont{JuliaMono}

\newcommand\calV{\mathcal{V}}
\newcommand\bbD{\mathbb{D}}
\newcommand\bbI{\mathbb{I}}
\newcommand\bbF{\mathbb{F}}
\newcommand\JEq{\mathtt{JEq}}
\newcommand\FTy{\mathtt{FTy}}
\newcommand\comp{\mathtt{comp}}
\newcommand\fst{\mathtt{fst}}
\newcommand\snd{\mathtt{snd}}
\newcommand\Ty{\mathtt{Ty}}
\newcommand\Ter{\mathtt{Ter}}
\newcommand\Path{\mathtt{Path}}
\newcommand\dM{\mathsf{dM}}
\newcommand\ttp{\mathtt{p}}
\newcommand\ttq{\mathtt{q}}
\newcommand\bfone{\mathbf{1}}
\newcommand\Set{\mathbf{Set}}
\newcommand\DM{\mathbf{DM}}
\newcommand\Id{\mathtt{Id}}
\newcommand\refl{\mathtt{refl}}
\newcommand\C{\mathbf{C}}
\newcommand\psC{\widehat{\mathbf{C}}}
\newcommand\extop{{\scalebox{0.4}[1]{\_}}.{\scalebox{0.8}[1]{\_}}}
\newcommand{\shortminus}{\raisebox{0.3ex}{\scalebox{0.8}[1.0]{-}}}
\newcommand{\mathshortminus}{\mathbin{\text{\normalfont\shortminus}}}
\newcommand{\ul}{{\scalebox{0.4}[1.0]{\_}}}

\DeclareMathOperator\FreeDeMorgan{FreeDeMorgan}
\DeclareMathOperator\obj{obj}
\DeclareMathOperator\y{\mathbf{y}}

\usepackage{unicode-math}
\setmathfont{STIX Two Math} 
%\setmathfont{Latin Modern Math}
%\setmathfont{XITS Math}[range=cal]
%\setmathfont{Euler Math}[range=cal]

\DeclareMathAlphabet{\mathcal}{OMS}{cmsy}{m}{n}

\newcommand\cubicalsets{\widehat{□}}

%----------------------------------------------------------------------------------------
%    TITLE PAGE
%----------------------------------------------------------------------------------------

\title{扩展类型(extension type)}
%\subtitle{Subtitle}

%\author{Pin-Yen Huang}

%\institute
%{
%    Department of Computer Science and Information Engineering \\
%    National Taiwan University % Your institution for the title page
%}
%\date{\today} % Date, can be changed to a custom date
\date{}

%----------------------------------------------------------------------------------------
%    PRESENTATION SLIDES
%----------------------------------------------------------------------------------------

\begin{document}

\begin{frame}
    % Print the title page as the first slide
    \titlepage
\end{frame}

\begin{frame}{背景}
\begin{itemize}
  \item $I \mid \Phi \subseteq \Psi$
  \item $I \mid \Psi \mid \Gamma \vdash A$
  \item $t : I \mid \Phi \mid \Gamma \vdash a(t) : A(t)$
  \item $t : I \mid \Psi \mid \Gamma \vdash b(t) : A(t)$, 满足 $b|_{\Phi}\equiv a$
  \item 上面的 $b$ 就是扩展类型 $\left\langle \prod_{t:\Psi} A(t) \middle|_{a}^{\Phi}\right\rangle$ 的一个项(term)
  \item 如果 $A$ 是一个常量类型, 扩展类型可简写为 $\left\langle\Psi \to A\middle|_{a}^{\Phi}\right\rangle$
  \item 类比从类型到类型的函数类型, 我们可以如下定义从形状(shape)到类型的函数类型: 可以取 $\Phi:\equiv\bot$, 我们有 $\mathtt{rec}_{\bot}$. 从 $\Psi$ 到 $A$ 的函数类型为 $\Psi \to A:\equiv\left\langle\Psi\to A\middle|_{\mathtt{rec}_{\bot}}^{\bot}\right\rangle$
\end{itemize}
\end{frame}

\begin{frame}[fragile,shrink]{rzk 中的 extension type}
rzk 中有两个不同的概念, 结合起来可以支持扩展类型
\begin{itemize}
  \item rzk 允许依赖于 cube 或 shape 的函数, 就是上面讨论的 $\Psi \to A:\equiv\left\langle\Psi\to A\middle|_{\mathtt{rec}_{\bot}}^{\bot}\right\rangle$
  \item 任何类型都允许“细化”(refinement), 也就是指定在 tope 限制下的值: 比如 $A\ [𝜙 ↦  x,\ 𝜓 ↦  y]$, 当 $𝜙$ 成立时, 求值为 $x$, 当 $𝜓$ 成立时, 求值为 $y$, 当然有一个前提, $x$ 和 $y$ 在 $(𝜙 \land 𝜓)$ 成立时要一样. $A$ 和 $A\ [\bot ↦ \mathtt{recBOT}]$ 是等价的.
\end{itemize}

例:\[\hom_A(a,b) :\equiv \left\langle\Delta^{1} \to A\middle|_{[a,b]}^{\partial\Delta^1}\right\rangle\]
在 rzk 中定义为

\begin{verbatim}
    #def hom (A : U) (a b : A) : U
        := (t : Δ¹) → A [t ≡ 0₂ ↦ a , t ≡ 1₂ ↦ b]
\end{verbatim}

\end{frame}


\end{document}